This is a reading aid, not a complete grammar. Names such as \code{target},
\code{pattern}, and \code{body} are placeholders; ellipses stand for omitted
parts. The linked sections explain evaluation, matching, prerequisites,
and the limits of the supplied implementation. Lisp notation and control
forms are introduced in Section~\ref{sec:lisp-reading}.

\paragraph{Objects, state, and matching.}
\begin{longtable}{L{0.39\textwidth}L{0.39\textwidth}L{0.12\textwidth}}
\textbf{Construct} & \textbf{Meaning} & \textbf{Section}\\\hline
\endhead
\code{[object name (script ...)]} & Create an object and bind its name;
omit the name for an anonymous object. & \ref{sec:objects}\\
\code{(state [n := 0])} & Object-definition entry for persistent state. & \ref{sec:state}\\
\code{(temporary [x := 0])} & Handler-local binding, not persistent state. & \ref{sec:objects}\\
\code{[x := expression]} & Change a binding. & \ref{sec:state}\\
\code{(=> pattern body)} & A script handling a matching message. & \ref{sec:objects}\\
\code{(=> pattern @ R from S body)} & Additionally bind reply destination
and sender; either header field can be omitted. & \ref{sec:objects}\\
\code{[head . tail]} & Pattern binding the head and entire tail of a pair;
dot notation is not an extra list element. & \ref{sec:lisp-reading}\\
\code{(match value (is pattern body) ...)} & Match a value inside a script;
\code{otherwise} supplies a fallback clause. & \ref{sec:objects}\\
\code{(=> pattern where test body)} & Guarded script selection, requiring
compatible heavy-meta machinery; not an ordinary \code{match} constraint. & \ref{sec:guards}\\
\end{longtable}

\paragraph{Sending, replying, and continuing.}
\begin{longtable}{L{0.39\textwidth}L{0.39\textwidth}L{0.12\textwidth}}
\textbf{Construct} & \textbf{Meaning} & \textbf{Section}\\\hline
\endhead
\code{[target <= message]} & Past-type send; do not wait for a reply. & \ref{sec:messages}\\
\code{[target <== message]} & Now-type send; wait for a reply. & \ref{sec:messages}\\
\code{[target <= message @ R]} & Past-type send with a specified reply destination. & \ref{sec:messages}\\
\code{[target <= message \$ f]} & Future-type send; route replies through
the loaded future library. & \ref{sec:futures}\\
\code{!expression} & Send a reply for the current message; not a return
from the rest of the handler. & \ref{sec:objects}\\
\code{[cont pattern body]} & Create a short, non-reifying continuation object. & \ref{sec:messages}\\
\code{[target <= message :express]} & Express send; requires both sender
and receiver machinery, not just this extra parameter. & \ref{sec:message-extensions}\\
\code{(atomic body ...)} & Receiver-local protection in the express
protocol; not a multi-object transaction. & \ref{sec:message-extensions}\\
\end{longtable}

\noindent\begin{minipage}{\linewidth}
\paragraph{Groups, identities, and reflection.}
\begin{longtable}{L{0.39\textwidth}L{0.39\textwidth}L{0.12\textwidth}}
\textbf{Construct} & \textbf{Meaning} & \textbf{Section}\\\hline
\endhead
\code{[group name ...]} & Create a group manager with optional custom
kernel objects. & \ref{sec:reflection}\\
\code{(group manager)} & Object-definition parameter selecting its
creation group. & \ref{sec:reflection}\\
\code{(meta-gen generator)} & Object/group definition parameter selecting
a metaobject generator. & \ref{sec:reflection}\\
\code{(evaluator evaluator-object)} & Group-definition parameter selecting
its evaluator. & \ref{sec:reflection}\\
\code{[meta object]} & Obtain its metaobject. & \ref{sec:reflection}\\
\code{[den metaobject]} & Obtain its lower-level denotation, not a copy. & \ref{sec:reflection}\\
\code{[group-of object]} & Query its group manager; custom generators must
provide the corresponding identity metadata. & \ref{sec:original-inherit}\\
\code{Me} / \code{group} & Current object denoted by \code{Id} / group
manager in \code{Gid}, not necessarily base-level application objects. & \ref{sec:reflection}\\
\code{(reify-args tag forms ...)} & Evaluate the tag, then send compiled
argument expressions and context to the current evaluator. & \ref{sec:reflection}\\
\code{(reify-args-to-meta tag forms ...)} & Send them to the current
object's metaobject instead. & \ref{sec:reflection}\\
\end{longtable}
\end{minipage}

\code{reify-args} and \code{reify-args-to-meta} are unavailable in a
non-reifying object's script. Extra send/object/assignment parameters
are interpreted by the selected machinery, not universally by their
spelling. Section~\ref{sec:reflection} explains the evaluation context;
Section~\ref{sec:files} separately collects loading and compilation commands.
